\section{并发更新与原子操作}
\label{l13:sec:atomic}

\subsection{读—改—写与更新丢失}

共享计数器的更新通常包含读取旧值、计算新值和写回三个步骤，称为读—改—写（read-modify-write）。例如，\mbox{\code{++bins[b]}} 在逻辑上执行“读出第 $b$ 个计数器，将它加一，再写回”。即使源程序只写了一条语句，也不能据此认为其他线程无法插入执行。

设共享计数器初值为 $0$，线程 $A$ 与线程 $B$ 各执行一次加一。每个线程分别持有局部变量 $r_A$、$r_B$；这些变量只记录该线程读取到的旧值。若 $A$ 完成整个更新后 $B$ 才读取，则 $A$ 读到 $0$、写入 $1$，$B$ 随后读到 $1$、写入 $2$。交换两个线程的先后顺序，最终也得到 $2$。然而，表\ref{l13:tab:lost-update}中的交错执行会丢失一次更新。

\begin{table}[htbp]
\centering
\caption{两个线程各自完成加一，最终计数却只有一}
\label{l13:tab:lost-update}
\begin{tabularx}{\textwidth}{cYYc}
\toprule
时刻 & 线程 $A$ & 线程 $B$ & 共享值 \\
\midrule
1 & 读取旧值，$r_A=0$ & & $0$ \\
2 & 计算 $r_A+1=1$ & & $0$ \\
3 & & 读取旧值，$r_B=0$ & $0$ \\
4 & 写回 $1$ & & $1$ \\
5 & & 计算 $r_B+1=1$ & $1$ \\
6 & & 写回 $1$ & $1$ \\
\bottomrule
\end{tabularx}
\end{table}

两个线程都完成了自己的指令，但它们根据同一个旧值计算，第二次写回覆盖了第一次更新。问题出在共享状态的访问没有形成正确的先后关系。多个线程并发访问同一内存位置、至少一次为写入，并且这些冲突访问缺少适当的同步或原子保证时，就存在数据竞争（data race）。上表只展示一种直观的错误交错；含数据竞争的程序不能依靠这几种手工列举的结果来保证行为。

类似的问题也出现在状态占用中。若用 $0$ 表示座位空闲、$1$ 表示已占用，两个请求都可能先读到 $0$，再分别写入 $1$，于是它们都认为自己取得了同一个座位。计数需要保护“读取后累加”的整体，占用操作需要保护“检查条件后修改”的整体。两者都要求把依赖旧状态的判断或计算与相应写回联系起来。

是否发生冲突由最终访问的地址决定。两个线程读取不同输入，并不保证它们写入不同输出。例如，不同键可能散列到同一个表项，不同顶点也可能向共同邻居累加贡献。直方图正具有这种结构：输入位置可以独立划分，而输入的值可能把多个线程映射到同一个计数器。

\subsection{互斥、线程束与执行进度}

互斥（mutual exclusion）要求同一时刻至多有一个参与者执行受保护的更新。将读取、计算和写回放入临界区（critical section），并在进入前获取锁、退出后释放锁，可以使整个更新按某种顺序完成。例如，对同一把互斥锁，只有 $A$ 释放后 $B$ 才能进入，或者只有 $B$ 释放后 $A$ 才能进入。两种顺序都避免了上面的覆盖问题。

这种方法依赖一个进度条件：获得锁的线程必须能够继续执行并最终释放锁。在传统的线程束（warp）共同推进模型中，同一束中的一个线程可能获得锁，其余线程则停留在等待锁的循环中。若执行调度一直停留在等待分支，使持锁线程的临界区与释放语句无法得到执行，就出现循环等待：等待者需要锁被释放，持锁者又无法越过当前执行路径去释放锁。这是把简单自旋锁直接搬入这种执行模型时可能发生死锁的原因。

上述死锁分析以线程束内的具体调度模型为前提。独立线程调度（independent thread scheduling）改变了线程束内线程推进的方式，因此分析锁时仍需同时检查参与线程的执行方式、内存可见性以及持锁者的进度。对于单个计数器加一这样的操作，可以直接使用硬件支持的原子更新，避免手工管理锁住和释放的临界区。

\subsection{原子更新的保证与边界}

原子操作（atomic operation）将指定地址上的读—改—写作为一个不可分割的操作完成。对于同一地址上、作用域相互覆盖的冲突原子操作，每次更新都必须看到某个完整的先前状态；硬件不会让另一原子更新插入这次操作的中间阶段。因此，同地址的冲突更新在效果上形成一个串行顺序。

若计数器初值为 $h_0$，所有 $m$ 次更新都通过原子加一完成，而且没有溢出或其他未同步访问，那么最终值就是 $h_0+m$。执行顺序可以未知，结果仍然确定，因为整数加法满足交换律和结合律。对于其他更新，顺序仍可能影响结果。例如，对同一位置分别执行原子交换，写入 $3$ 和写入 $7$，最终值取决于哪个交换最后完成。

原子保证只覆盖该操作指定的状态转移。对两个不同计数器先后执行两次原子加法，并不会自动把二者合成一个不可分割的事务；另一个线程可能观察到其中一次已经完成、另一次尚未完成。对单个计数器的原子加法也不提供“所有线程都已完成”的通知。后面在共享直方图中使用的块内同步，负责保证一个阶段全部结束后才开始下一阶段；原子加法则负责该阶段内每次冲突更新的正确性。

\begin{keybox}[title={原子性与顺序}]
同一地址的原子更新不会互相覆盖掉未完成的中间结果，但线程之间谁先更新通常没有预定顺序。计数结果的确定性还依赖更新运算本身；完整结果应在所有更新结束后读取。
\end{keybox}

\subsection{CUDA 原子接口与比较交换}

CUDA 提供由编译器识别的内建原子函数。\mbox{\code{atomicAdd}} 接受目标地址和增量两个参数，原子地读出旧值、加上增量、写回结果，并返回更新前的旧值。例如，一次调用向计数器加入 $3$，若原子动作读到 $4$，它就写入 $7$，并向调用者返回 $4$。多个线程争用同一地址时，返回值也服从实际形成的串行顺序。

\mbox{\code{atomicAdd}} 有不同数据类型的重载，常见类型包括整数与浮点数；具体重载的支持取决于设备和数据类型。计数问题通常使用整数。其他原子原语包括减法、最小值、最大值、按位与、按位或、按位异或、交换以及比较交换。对于可交换的整数累加，原子操作保持计数正确；浮点加法还受到舍入和累加顺序的影响。

比较交换（compare-and-swap，CAS）把“检验当前位置是否仍为预期值”与“条件成立时写入新值”合成一个原子动作。用 $v$ 表示操作开始时该地址的值，$e$ 表示预期值，$d$ 表示希望写入的值，其状态变化为
\begin{equation}
v_{\mathrm{after}}=
\begin{cases}
d, & v=e,\\
v, & v\ne e.
\end{cases}
\label{l13:eq:cas}
\end{equation}
在抽象描述中，可以让 CAS 返回“成功”或“失败”。CUDA 的 \mbox{\code{atomicCAS}} 接口依次接受目标地址、预期值和目标新值三个参数。无论替换是否成功，它都返回操作之前观察到的旧值。因此，返回旧值等于预期值，才是替换成功的判断条件。

这一区别会直接影响程序逻辑。如果旧值为 $0$、预期值也为 $0$，替换成功时返回 $0$；直接将返回值当作布尔量，会把这次成功当成失败。反过来，如果当前值为 $5$、预期值为 $0$，替换失败时返回 $5$；将其转换为布尔量又会得到真。普通的条件判断加赋值本身也无法实现式\eqref{l13:eq:cas}要求的不可分割性，真正的保证来自原子原语。

对座位占用状态，可以执行 \mbox{\code{old = atomicCAS(&seat, 0u, 1u)}}。若返回 $0$，说明本次调用取得了此前空闲的座位；若返回 $1$，说明调用时座位已被占用。两个并发请求都使用这一原语时，只有一个能够把同一次空闲状态成功改为占用状态。

无后缀的传统 CUDA 原子接口采用设备作用域；对于全局计数器，这允许同一设备中不同线程块共同原子更新。它们不会自动建立其他地址的读写顺序。共享内存对象则属于各自线程块，只有该块中的线程参与访问。地址必须有效并满足类型对齐要求，初始化需要先于并发更新，所有冲突访问也必须遵守相容的原子与同步规则。

\subsection{用比较交换构造原子加法}

CAS 可以构造“根据旧值计算新值”的更新。线程先取得一个估计值，计算目标值，然后尝试仅在共享值仍与估计相等时写入。如果比较时共享值与估计值不同，CAS 会失败并返回新的观测值；当前线程据此重新计算和重试。这样，线程私有的计算可以被打断，最终提交更新的动作仍然是原子的。

代码\ref{l13:lst:cas-add}用无符号整数实现这一过程。函数名与 CUDA 内建的 \mbox{\code{atomicAdd}} 区分开，且整个实现中的共享访问都使用 CAS。初始调用 \mbox{\code{atomicCAS(address, 0u, 0u)}} 只用于取得当前值：当前值为 $0$ 时写回相同的 $0$，否则保持原值。

\noindent\begin{minipage}{\textwidth}
\begin{lstlisting}[caption={由 CAS 构造无符号整数原子加法},label={l13:lst:cas-add}]
__device__ unsigned int atomic_add_cas(
    unsigned int* address, unsigned int value)
{
    unsigned int old = atomicCAS(address, 0u, 0u);
    while (true) {
        unsigned int assumed = old;
        old = atomicCAS(address, assumed, assumed + value);
        if (old == assumed) {
            return assumed;
        }
    }
}
\end{lstlisting}
\end{minipage}

变量 \mbox{\code{assumed}} 保存本轮准备提交时的预期旧值；\mbox{\code{old}} 在 CAS 返回后变为该原子动作实际观察到的值。若二者相同，说明比较时共享值确实等于 \mbox{\code{assumed}}，写入的正是这个当前值加上 \mbox{\code{value}}，函数返回此次更新前的值。若二者不同，本次 CAS 没有修改共享状态，下一轮直接使用返回值作为新估计。

例如，初始计数为 $0$，$A$、$B$ 各加 $1$，两者都先估计为 $0$。$A$ 的 CAS 将 $0$ 改成 $1$，返回 $0$，因此成功。$B$ 随后尝试把 $0$ 改成 $1$，但当前值已经是 $1$，于是 CAS 返回 $1$ 且不写入。$B$ 下一轮改为尝试把 $1$ 改成 $2$，成功后返回 $1$。最终计数为 $2$。

正确性的关键是成功 CAS 的时刻。在线程准备估计值或计算加法期间，其他线程都可以继续更新；只有当提交时的共享值满足比较条件，当前计算才会生效。因此，函数整体可以被看作在成功 CAS 时刻完成了一次加法。竞争越激烈，失败重试可能越多，不能为每个线程保证固定的重试次数或取得更新机会的顺序。已有合适的内建原子加法时，应直接使用内建函数；这个构造用于解释 CAS 如何提供更一般的更新机制。

无符号整数在溢出时按其位宽回绕。若把它作为通常意义上的计数器，仍须保证实际计数不超过可表示范围。原子性保证更新不会相互丢失，计数容量则决定最终数值能否完整表示真实数量。
